\chapter{Données et méthodologie}
\label{chapitre:methodologie}

Ce chapitre présente les fondations empiriques et méthodologiques du système. Il décrit
d'abord le corpus de données capteurs sur lequel repose l'ensemble du travail, puis les
transformations appliquées aux signaux bruts pour en extraire des variables exploitables,
l'architecture du pipeline de détection, et enfin les métriques et protocoles servant à
évaluer le système. L'objectif est de rendre explicite chaque choix méthodologique, depuis
l'ingestion des données jusqu'aux critères de validation.

% =====================================================================
\section{Corpus de données capteurs}
\label{sec:donnees}
% =====================================================================

Le corpus principal du projet, c'est-à-dire l'ensemble des mesures capteurs exploitées dans
ce mémoire, est constitué de données issues de capteurs IceQube (IceRobotics Ltd),
collectées à la ferme laitière du campus Macdonald de l'Université McGill
(Sainte-Anne-de-Bellevue, Québec). Les capteurs, des
accéléromètres tri-axiaux fixés à la patte arrière des animaux
(Section~\ref{sec:technologies_capteurs}), aggrègent leurs mesures par tranches de quinze
minutes, un format courant dans les études de suivi bovin
\citep{ledgerwood2010lying,benaissa2019calving}. Les données brutes sont exportées au
format CSV depuis le logiciel du capteur. Elles sont également consultables via la
plateforme web CowAlert (IceRobotics), qui fournit une interface de visualisation
en ligne des mêmes mesures. La qualité de ces données conditionnant directement la validité
des analyses, leur intégrité a fait l'objet de vérifications préliminaires avant toute
exploitation.

\subsection{Vérification de l'intégrité des données}
Avant toute analyse, deux vérifications préliminaires ont été conduites afin de s'assurer
de la fiabilité du corpus.

\textit{Cohérence entre sources.}
Les données brutes exportées depuis le capteur IceQube ont été comparées à celles affichées
sur la plateforme CowAlert. Sur 752~observations appariées (rééchantillonnées à 1~heure),
les corrélations obtenues sont quasi parfaites pour les deux variables d'activité
principales, le nombre de pas (\textit{Steps}) et l'indice de mouvement
(\textit{Motion Index}), toutes deux définies au Tableau~\ref{tab:variables_brutes}:
$r_{\text{Pearson}} = 1.000$ pour \textit{Steps} et \textit{Motion Index}, $r_{\text{Spearman}} = 1.000$
et $0.999$ respectivement. Cette correspondance
confirme que CowAlert restitue fidèlement les données brutes du capteur et que les fichiers
CSV utilisés dans ce projet reflètent exactement les mesures enregistrées.

\textit{Fuseau horaire.}
Une analyse des horodatages a révélé que les données exportées par le capteur IceQube sont
enregistrées en temps universel coordonné (UTC), avec un décalage de $-5$~heures par rapport
à l'heure locale de la ferme (fuseau EST, Montréal). Ainsi, un bin horodaté
\texttt{12:00 UTC} correspond à \texttt{07:00} heure locale. Ce décalage n'affecte pas le pipeline de détection, qui opère sur des fenêtres glissantes
relatives indépendantes du fuseau, mais il a été pris en compte dans l'interprétation des
profils circadiens et des visualisations.

La version de référence contient
37\,839 lignes couvrant 28~vaches, sur la période du 1\textsuperscript{er} octobre~2023
au 11~novembre~2023. Les variables comportementales disponibles sont résumées dans le
Tableau~\ref{tab:dataset}.

\begin{table}[H]
    \centering
    \caption{Description du corpus de données capteurs.}
    \label{tab:dataset}
    \begin{tabular}{ll}
        \toprule
        \textbf{Propriété}         & \textbf{Valeur} \\
        \midrule
        Nombre de lignes           & 37\,839 \\
        Nombre de vaches           & 28 \\
        Période                    & 1\textsuperscript{er} oct. -- 11 nov. 2023 \\
        Intervalle temporel        & 15 minutes \\
        Lignes médianes par vache  & $\approx$ 564 \\
        Lignes moyennes par vache  & $\approx$ 1\,350 \\
        \bottomrule
    \end{tabular}
\end{table}

Le Tableau~\ref{tab:stats_descriptives} présente les statistiques descriptives des
principales variables comportementales, calculées sur l'ensemble du corpus
(28~vaches, 37\,839~lignes en bins de 15~minutes) après filtrage des bins à couverture
insuffisante.

\begin{table}[H]
    \centering
    \caption{Statistiques descriptives des variables comportementales brutes.}
    \label{tab:stats_descriptives}
    \begin{tabular}{lcccccc}
        \toprule
        \textbf{Variable} & \textbf{Médiane} & \textbf{Moy.} & \textbf{Éc.-type}
            & \textbf{Q1} & \textbf{Q3} & \textbf{Max} \\
        \midrule
        \textit{Steps}              & 0     & 4.0   & 15.4  & 0     & 3     & 387 \\
        \textit{Motion Index}       & 5.0   & 25.5  & 103.7 & 0.0   & 19.0  & 3351 \\
        \textit{Lying Time} (min)   & 15.0  & 9.3   & 6.9   & 0.0   & 15.0  & 15.0 \\
        \textit{Standing Time} (min)& 0.0   & 5.7   & 6.9   & 0.0   & 15.0  & 15.0 \\
        \textit{Transitions}        & 0     & 0.2   & 0.4   & 0     & 0     & 6 \\
        \bottomrule
    \end{tabular}
\end{table}

Trois caractéristiques du corpus méritent d'être soulignées. D'abord, les distributions
sont fortement asymétriques: les médianes de \textit{Steps} et \textit{Motion Index} sont nettement
inférieures aux moyennes, ce qui indique la présence de valeurs extrêmes. Ensuite, la
variable \textit{Lying Time} est bornée à 15~minutes (la durée du bin), ce qui crée un effet plafond
naturel. Enfin, les transitions posturales sont rares par bin (médiane de 0), ce qui rend
cette variable plus informative en agrégation glissante qu'en valeur ponctuelle.

La Figure~\ref{fig:distribution_steps} illustre l'asymétrie de la distribution des \textit{Steps}
sur l'ensemble du corpus (37\,839~bins, médiane~=~0, moyenne~=~4)~: plus de 90\,\% des bins
de 15~minutes comptent moins de 10~pas, l'animal étant le plus souvent au repos sur un
intervalle aussi court. Cette non-normalité
justifie l'utilisation de statistiques robustes (médiane, MAD) plutôt que de la moyenne
et de l'écart-type pour le calcul des z-scores \citep{leys2013detecting}.

\begin{figure}[H]
    \centering
    \begin{tikzpicture}
        \begin{axis}[
            ybar interval,
            width=11cm, height=6cm,
            ylabel={Fréquence relative (\%)},
            xlabel={Steps par bin de 15 minutes},
            xmin=0, xmax=100,
            ymin=0, ymax=100,
            ymajorgrids=true,
            grid style={black!8},
            axis line style={black!60},
            tick style={black!40},
            x tick label style={font=\scriptsize},
            area style,
        ]
            \addplot[fill=ifblue!25, draw=ifblue!50] coordinates {
                (0,93.4) (10,4.1) (20,0.8) (30,0.5) (40,0.3)
                (50,0.2) (60,0.2) (70,0.2) (80,0.2) (90,0.2) (100,0)
            };
            % Ligne moyenne (dans le repère de l'axe)
            \draw[loforange, thick, densely dotted] (axis cs:4,0) -- (axis cs:4,95);
            % Labels dans une boîte en haut à droite pour éviter toute superposition
            \node[font=\scriptsize, anchor=north east, fill=white, fill opacity=0.85,
                  text opacity=1, inner sep=4pt, rounded corners=2pt, draw=black!15,
                  align=left]
                at (axis cs:98,92) {%
                \textcolor{rawred}{médiane = 0}\\
                \textcolor{loforange}{$\cdots$\; moyenne = 4}\\
                {\footnotesize tronqué à 100 (max~=~387)}%
            };
        \end{axis}
    \end{tikzpicture}
    \caption{Distribution des \textit{Steps} sur l'ensemble du corpus.}
    \label{fig:distribution_steps}
\end{figure}

La Figure~\ref{fig:couverture_corpus} illustre la couverture temporelle du corpus par
vache, montrant que le volume de données varie fortement selon les animaux. La
distribution est nettement bimodale~: une partie des vaches dispose d'environ
564~bins ($\approx$6~jours), une autre d'environ 2\,450~bins ($\approx$25~jours), cet
écart reflétant des durées de port du capteur différentes plutôt que des trous dans la
mesure. La médiane est de 564~bins par vache et la moyenne d'environ 1\,350~bins
($\approx$14~jours)~; au sein de la fenêtre d'observation propre à chaque animal, la
couverture par bin reste quasi complète.

\begin{figure}[H]
    \centering
    \begin{tikzpicture}
        \begin{axis}[
            ybar,
            bar width=8pt,
            width=12cm, height=5.5cm,
            ylabel={Nombre de bins},
            xlabel={Vaches (triées par volume de données)},
            ymin=0, ymax=2700,
            ymajorgrids=true,
            grid style={black!8},
            xtick=\empty,
            xmin=0, xmax=29,
            axis line style={black!60},
            tick style={black!40},
            nodes near coords=\empty,
        ]
            \addplot[fill=ifblue!25, draw=ifblue!50] coordinates {
                (1,244) (2,412) (3,441) (4,564) (5,564)
                (6,564) (7,564) (8,564) (9,564) (10,564)
                (11,564) (12,564) (13,564) (14,564) (15,564)
                (16,564) (17,2444) (18,2444) (19,2449) (20,2449)
                (21,2449) (22,2449) (23,2450) (24,2450) (25,2450)
                (26,2450) (27,2452) (28,2474)
            };
            % Ligne médiane
            \draw[rawred, dashed, thick] (axis cs:0,564) -- (axis cs:29,564);
            % Ligne moyenne
            \draw[loforange, densely dotted, thick] (axis cs:0,1350) -- (axis cs:29,1350);
            % Légende en haut à gauche (zone vide)
            \node[font=\scriptsize, anchor=north west, fill=white, fill opacity=0.85,
                  text opacity=1, inner sep=4pt, rounded corners=2pt, draw=black!15,
                  align=left]
                at (axis cs:1,2600) {%
                \textcolor{rawred}{- - -\; médiane = 564 bins}\\
                \textcolor{loforange}{$\cdots$\; moyenne $\approx$ 1\,350 bins}%
            };
        \end{axis}
    \end{tikzpicture}
    \caption{Distribution du nombre de bins par vache dans le corpus.}
    \label{fig:couverture_corpus}
\end{figure}

La Figure~\ref{fig:heatmap_couverture} illustre schématiquement le principe de suivi de
la couverture par vache et par semaine. Les pourcentages
indiquent la proportion de bins valides par semaine~; les intensités de couleur
reflètent le nombre de bins valides, révélant les périodes de maintenance capteur ou de
couverture partielle. Sur le corpus réel, la couverture par bin reste quasi complète à
l'intérieur de la fenêtre d'observation propre à chaque vache, de sorte qu'aucune vache
n'atteint le seuil de filtrage (\texttt{coverage\_min\_pct}~=~25\%), ce qui confirme la
bonne qualité générale du corpus~; les variations de volume entre animaux
(Figure~\ref{fig:couverture_corpus}) proviennent surtout de durées de port différentes.

\begin{figure}[H]
    \centering
    \begin{tikzpicture}[
        cell/.style={minimum width=1.2cm, minimum height=0.35cm, font=\tiny, inner sep=0pt},
    ]
        % Titre colonnes (semaines)
        \foreach \j/\lab in {1/S1, 2/S2, 3/S3, 4/S4, 5/S5, 6/S6} {
            \node[font=\scriptsize\bfseries, black] at (\j*1.2, 0.4) {\lab};
        }
        % Titre lignes (vaches - échantillon)
        \foreach \i/\lab in {1/V1, 2/V5, 3/V10, 4/V15, 5/V20, 6/V25, 7/V28} {
            \node[font=\scriptsize, black!60, anchor=east] at (0.3, -\i*0.45) {\lab};
        }
        % Cellules heatmap (intensité = couverture)
        % V1
        \foreach \j/\c in {1/95, 2/92, 3/88, 4/90, 5/85, 6/78} {
            \pgfmathsetmacro{\intensity}{\c/100*40}
            \fill[ifblue!\intensity] (\j*1.2-0.5, -0.45*1+0.15) rectangle (\j*1.2+0.5, -0.45*1-0.15);
            \node[cell] at (\j*1.2, -0.45*1) {\c\%};
        }
        % V5
        \foreach \j/\c in {1/98, 2/97, 3/95, 4/96, 5/94, 6/92} {
            \pgfmathsetmacro{\intensity}{\c/100*40}
            \fill[ifblue!\intensity] (\j*1.2-0.5, -0.45*2+0.15) rectangle (\j*1.2+0.5, -0.45*2-0.15);
            \node[cell] at (\j*1.2, -0.45*2) {\c\%};
        }
        % V10
        \foreach \j/\c in {1/92, 2/90, 3/85, 4/88, 5/82, 6/75} {
            \pgfmathsetmacro{\intensity}{\c/100*40}
            \fill[ifblue!\intensity] (\j*1.2-0.5, -0.45*3+0.15) rectangle (\j*1.2+0.5, -0.45*3-0.15);
            \node[cell] at (\j*1.2, -0.45*3) {\c\%};
        }
        % V15
        \foreach \j/\c in {1/88, 2/85, 3/80, 4/82, 5/78, 6/70} {
            \pgfmathsetmacro{\intensity}{\c/100*40}
            \fill[ifblue!\intensity] (\j*1.2-0.5, -0.45*4+0.15) rectangle (\j*1.2+0.5, -0.45*4-0.15);
            \node[cell] at (\j*1.2, -0.45*4) {\c\%};
        }
        % V20
        \foreach \j/\c in {1/96, 2/94, 3/92, 4/93, 5/90, 6/88} {
            \pgfmathsetmacro{\intensity}{\c/100*40}
            \fill[ifblue!\intensity] (\j*1.2-0.5, -0.45*5+0.15) rectangle (\j*1.2+0.5, -0.45*5-0.15);
            \node[cell] at (\j*1.2, -0.45*5) {\c\%};
        }
        % V25
        \foreach \j/\c in {1/90, 2/88, 3/83, 4/85, 5/80, 6/72} {
            \pgfmathsetmacro{\intensity}{\c/100*40}
            \fill[ifblue!\intensity] (\j*1.2-0.5, -0.45*6+0.15) rectangle (\j*1.2+0.5, -0.45*6-0.15);
            \node[cell] at (\j*1.2, -0.45*6) {\c\%};
        }
        % V28
        \foreach \j/\c in {1/97, 2/95, 3/93, 4/94, 5/92, 6/90} {
            \pgfmathsetmacro{\intensity}{\c/100*40}
            \fill[ifblue!\intensity] (\j*1.2-0.5, -0.45*7+0.15) rectangle (\j*1.2+0.5, -0.45*7-0.15);
            \node[cell] at (\j*1.2, -0.45*7) {\c\%};
        }
        % Légende
        \node[font=\tiny, black!50, anchor=west] at (8, -0.45*3) {Couverture};
        \fill[ifblue!12] (8, -0.45*4) rectangle (8.6, -0.45*4+0.25);
        \node[font=\tiny, black!50, anchor=west] at (8.8, -0.45*4+0.12) {70\%};
        \fill[ifblue!30] (8, -0.45*5) rectangle (8.6, -0.45*5+0.25);
        \node[font=\tiny, black!50, anchor=west] at (8.8, -0.45*5+0.12) {85\%};
        \fill[ifblue!40] (8, -0.45*6) rectangle (8.6, -0.45*6+0.25);
        \node[font=\tiny, black!50, anchor=west] at (8.8, -0.45*6+0.12) {95\%+};
    \end{tikzpicture}
    \caption{Illustration schématique du suivi de couverture temporelle hebdomadaire
    (échantillon de 7~vaches, oct.--nov.~2023).}
    \label{fig:heatmap_couverture}
\end{figure}

Chaque ligne associe un identifiant animal, un horodatage de début et de fin d'intervalle,
ainsi que sept variables comportementales directement exploitables, présentées dans le
Tableau~\ref{tab:variables_brutes}.

\begin{table}[H]
    \centering
    \caption{Variables comportementales brutes disponibles dans le corpus.}
    \label{tab:variables_brutes}
    \begin{tabular}{lll}
        \toprule
        \textbf{Variable}     & \textbf{Type}   & \textbf{Description} \\
        \midrule
        \textit{Steps}                 & Entier          & Nombre de pas cumulés sur l'intervalle \\
        \textit{Motion Index}          & Réel            & Indice d'activité motrice global \\
        \textit{Lying Time}            & Réel (min)      & Temps passé couché \\
        \textit{Standing Time}         & Réel (min)      & Temps passé debout \\
        \textit{Transitions}           & Entier          & Changements couché $\leftrightarrow$ debout \\
        \textit{Transitions Up}        & Entier          & Changements couché $\rightarrow$ debout \\
        \textit{Transitions Down}      & Entier          & Changements debout $\rightarrow$ couché \\
        \bottomrule
    \end{tabular}
\end{table}

Le choix d'utiliser des variables comportementales déjà agrégées est cohérent avec une
approche déployable sur capteurs d'élevage. La littérature sur la boiterie montre que les
signaux utiles sont fréquemment des descripteurs dérivés de l'activité, du couchage ou des
transitions, plutôt qu'un flux inertiel brut seul, car les effets de la boiterie sont souvent
indirects et fortement dépendants de l'individu
\citep{oleary2020accelerometers,alsaaod2012electronic,lavrova2023lameness}.

Les données capteurs présentent plusieurs contraintes structurelles. D'abord, aucune annotation
vétérinaire synchronisée par tranche horaire n'est disponible; le projet ne peut donc pas
être formulé comme une classification supervisée standard.
Ensuite, les séries sont hétérogènes entre animaux: certains profils sont plus actifs, plus
mobiles ou plus stables que d'autres. Enfin, la couverture capteur n'est pas constante, ce qui
impose une gestion explicite des trous et des débuts de série. Ces limites sont cohérentes
avec l'état de la littérature en détection automatisée de boiterie bovine
\citep{qiao2021review,merck_hoof_origin_lameness_cattle}.

\subsection{Qualité des données et critères de filtrage}

La qualité des données capteurs est un facteur déterminant pour la fiabilité du pipeline.
Quatre situations de non-qualité ont été identifiées et traitées lors du prétraitement:

\begin{enumerate}
    \item \textit{Couverture insuffisante}: certaines vaches présentent moins de 200~bins
    exploitables sur la période, ce qui ne permet pas d'estimer des statistiques de
    référence fiables. Ces animaux sont donc exclus de l'analyse. Le seuil de 200~bins
    (environ 3~jours de données continues) est cohérent avec la nécessité de calculer
    des z-scores (écarts normalisés par rapport à la médiane robuste de chaque vache) et
    des moyennes mobiles sur au moins 7~jours glissants;
    \item \textit{Valeurs aberrantes ponctuelles}: des valeurs de \textit{Steps} ou de \textit{Motion Index}
    ponctuellement très élevées (capteur perturbé, mouvement de manipulation) sont traitées
    par la normalisation robuste (médiane/IQR) plutôt que par un écrêtage explicite, ce qui
    évite d'introduire des seuils arbitraires;
    \item \textit{Trous temporels}: les interruptions dans la série temporelle (maintenance
    capteur, changement de batterie) sont gérées par la logique de fenêtre glissante. Cette
    logique ne requiert pas de continuité stricte, mais seulement un nombre minimum de bins
    valides dans chaque fenêtre;
    \item \textit{Homogénéité temporelle}: les bins de 15~minutes fournissent une granularité
    suffisante pour capturer les variations comportementales liées à la boiterie, tout en
    restant compatibles avec les fréquences d'échantillonnage standard des capteurs
    d'élevage \citep{ledgerwood2010lying}.
\end{enumerate}

\subsection{Variabilité inter-individuelle}

La variabilité entre animaux est un défi connu dans l'analyse de données capteurs en élevage
\citep{oleary2020accelerometers,nechanitzky2016analysis}. Les profils d'activité diffèrent
considérablement d'une vache à l'autre en raison de facteurs tels que la parité, le stade
de lactation, le rang social et le tempérament individuel.

Le Tableau~\ref{tab:variabilite_inter} illustre cette hétérogénéité sur les principales
variables du corpus. Les coefficients de variation (CV) inter-vaches montrent que certaines
variables (\textit{Transitions}, \textit{Motion Index}) sont plus dispersées que d'autres (\textit{Lying Time}),
ce qui justifie l'approche de modélisation par vache adoptée dans ce pipeline.

\begin{table}[H]
    \centering
    \caption{Variabilité inter-individuelle des principales variables comportementales.}
    \label{tab:variabilite_inter}
    \begin{tabular}{lccc}
        \toprule
        \textbf{Variable} & \textbf{Médiane inter-vaches}
        & \textbf{IQR inter-vaches} & \textbf{CV (\%)} \\
        \midrule
        \textit{Steps}              & 5.0          & 3.3--6.0       & 50 \\
        \textit{Motion Index}       & 27.1         & 22.2--37.6     & 40 \\
        \textit{Lying Time} (min)   & 8.4          & 7.0--9.5       & 24 \\
        \textit{Standing Time} (min)& 6.6          & 5.5--8.0       & 32 \\
        \textit{Transitions}        & 0.2          & 0.2--0.2       & 41 \\
        \bottomrule
    \end{tabular}
\end{table}

Une telle variabilité motive deux choix méthodologiques fondamentaux. D'abord,
l'entraînement du détecteur d'anomalies (Isolation Forest) séparément pour chaque vache,
afin que le comportement ``normal'' soit défini relativement à l'individu et non au
troupeau. Ensuite, l'utilisation de z-scores calculés par vache pour les règles métier,
de sorte qu'un même seuil absolu ait une signification comparable entre animaux aux
profils très différents. Ces choix sont détaillés dans les sections suivantes.

\subsection{Choix de la résolution temporelle}

Le choix de la résolution temporelle d'agrégation résulte d'un compromis entre la
granularité du signal et sa stabilité. Le Tableau~\ref{tab:choix_resolution} compare trois
résolutions candidates et leurs implications pour la détection de boiterie.

\begin{table}[H]
    \centering
    \caption{Comparaison des résolutions temporelles candidates pour l'agrégation capteur.}
    \label{tab:choix_resolution}
    \begin{tabular}{lccc}
        \toprule
        \textbf{Résolution} & \textbf{Bins/jour} & \textbf{Avantage} & \textbf{Inconvénient} \\
        \midrule
        5 minutes & 288 & Granularité fine & Bruit élevé, bins vides fréquents \\
        \textbf{15 minutes} & \textbf{96} & \textbf{Compromis} & \textbf{Retenu} \\
        1 heure & 24 & Signal lissé & Perte de précision temporelle \\
        \bottomrule
    \end{tabular}
\end{table}

De cette comparaison, la résolution de 15~minutes ressort comme le meilleur compromis.
Elle correspond en outre au pas natif du capteur utilisé dans ce projet et reste cohérente
avec les intervalles rapportés dans la littérature sur les capteurs d'activité bovine:
\cite{ledgerwood2010lying} recommandent un intervalle de 1 à 5~minutes pour le comportement
de couchage, tandis que \cite{alsaaod2012electronic} utilisent des agrégations horaires
pour la détection pédométrique. Notre choix de 15~minutes se situe
entre ces deux extrêmes, offrant une granularité suffisante pour capturer les transitions
comportementales tout en limitant le bruit des bins à faible couverture.

À cette résolution, chaque vache produit environ 96~bins par jour, soit $\approx$1\,350~bins
sur 14~jours. Ce volume est suffisant pour Isolation Forest, qui nécessite typiquement
quelques centaines de points pour estimer des chemins d'isolation
fiables \citep{liu2008isolation}.

% =====================================================================
\subsection{Corpus complémentaire: ferme du campus Macdonald (automne 2019)}
\label{subsec:corpus_mcgill}
% =====================================================================

En complément du corpus principal, un second jeu de données a été collecté à la même
exploitation, la ferme laitière du campus Macdonald de l'Université McGill
(Sainte-Anne-de-Bellevue, Québec), lors de l'automne~2019. Ce corpus est
constitué de données IceTag (IceRobotics Ltd), prédécesseur de l'IceQube utilisé en 2023,
produisant les mêmes variables comportementales dans le même format de 15~minutes.

\begin{table}[H]
    \centering
    \caption{Description du corpus complémentaire McGill IceTag (automne~2019).}
    \label{tab:dataset_mcgill}
    \begin{tabular}{ll}
        \toprule
        \textbf{Propriété}              & \textbf{Valeur} \\
        \midrule
        Nombre de bins (15~min)         & 93\,187 \\
        Nombre de vaches                & 30 \\
        Période                         & 11~nov.~2019 -- 14~déc.~2019 \\
        Durée par vache                 & 34~jours (32~pour 6~vaches) \\
        Bins médians par vache          & $\approx$ 3\,143 \\
        Variables disponibles           & Identiques au corpus principal \\
        Vaches avec score SLS           & 12 \\
        Vaches sans score SLS           & 18 \\
        Source des scores SLS           & Exercice d'évaluation hiver~2019\\
                                        & (baseline janv., mi-parcours mars~2019) \\
        \bottomrule
    \end{tabular}
\end{table}

\textit{Labels cliniques (SLS).}
Les 12~vaches ayant un score de locomotion standardisé (SLS, somme de quatre critères
posturaux, échelle 0--8) ont été évaluées en deux temps lors de l'hiver~2019. Le
Tableau~\ref{tab:mcgill_sls_groups} résume la distribution par groupe clinique~:
SLS~$\geq$~2 correspond au seuil d'alerte clinique retenu dans la littérature
\citep{sprecher1997lameness}, et les 18~vaches sans score sont conservées pour l'analyse
non supervisée sans utilisation des labels.

\begin{table}[H]
    \centering
    \small
    \caption{Distribution des 30~vaches du corpus McGill par groupe clinique.}
    \label{tab:mcgill_sls_groups}
    \begin{tabular}{lccc}
        \toprule
        \textbf{Groupe} & \textbf{N} &
        \textbf{SLS baseline (janv.~2019)} &
        \textbf{SLS mi-parcours (mars~2019)} \\
        \midrule
        Boiteuse persistante  & 1 & $\geq$2 & $\geq$2 \\
        Boiteuse aggravée     & 1 & $\geq$2 & $\geq$2 (progression) \\
        Boiteuse améliorée    & 2 & $\geq$2 & $<$2 \\
        Légère (SLS~=~1)      & 3 & 1       & 0--2 \\
        Saine et dégradée     & 5 & 0       & 0--1 \\
        \midrule
        \textit{Sans score SLS} & 18 & — & — \\
        \midrule
        \textbf{Total}        & \textbf{30} & & \\
        \bottomrule
    \end{tabular}
\end{table}

\textit{Rôle dans le projet.}
Le corpus complémentaire n'est pas utilisé dans la validation principale du pipeline, qui repose sur le
corpus IceQube 2023 (Section~\ref{sec:donnees},
Tableau~\ref{tab:dataset}). Il constitue un banc de test naturel pour deux questions
complémentaires:
\begin{enumerate}
    \item \textit{Généralisation intra-site}: le pipeline entraîné sur les données 2023
    produit-il des alertes cohérentes sur des vaches différentes de la même ferme,
    avec un capteur de la même famille?
    \item \textit{Concordance avec les labels cliniques}: les vaches classées
    \textit{boiteuses persistantes} ou \textit{aggravées} génèrent-elles
    significativement plus d'alertes que les vaches saines?
\end{enumerate}

\textit{Limites de l'alignement temporel.}
Les scores SLS ont été évalués en janvier et mars~2019, soit 8 à 9~mois \textit{avant}
la collecte des données IceTag (automne~2019). Le statut clinique d'une vache peut évoluer
significativement sur cet horizon. Cette désynchronisation est assumée: elle n'invalide pas
l'analyse, mais impose de l'interpréter comme une vérification de la cohérence
comportementale à long terme plutôt que comme une validation clinique synchrone.
La conversion des données brutes en format pipeline est réalisée par le script
\texttt{convert\_mcgill\_to\_pipeline.py} (dossier \texttt{mcgill\_iot\_cattle/}), qui
produit \texttt{mcgill\_brut.csv} (93\,187~lignes) selon le même schéma que
\texttt{data/brut.csv}.

\FloatBarrier

% =====================================================================
\section{Architecture du pipeline}
% =====================================================================

Le pipeline de détection se décompose en quatre étapes séquentielles, illustrées dans la
Figure~\ref{fig:pipeline}. Chaque étape est implémentée dans un module logiciel distinct,
ce qui permet de les tester et de les faire évoluer indépendamment.

\begin{figure}[H]
    \centering
    \begin{tikzpicture}[
        block/.style={rectangle, draw=black!70, rounded corners, minimum height=1cm,
                      minimum width=2.5cm, align=center, font=\footnotesize,
                      inner sep=10pt},
        arrow/.style={-{Stealth[length=2.5mm]}, thick, black!80},
    ]
        % Blocs
        \node[block, fill=gray!12] (io) {1. Chargement\\[-7pt]et filtrage};
        \node[block, fill=ifblue!15, right=0.5cm of io] (feat) {2. Extraction\\[-7pt]de \textit{features}};
        \node[block, fill=loforange!12, right=0.5cm of feat] (if) {3. Détection\\[-7pt](IF)};
        \node[block, fill=rulegreen!15, right=0.5cm of if] (rules) {4. Règles métier\\[-7pt]et notifications};

        % Fleches
        \draw[arrow] (io) -- (feat);
        \draw[arrow] (feat) -- (if);
        \draw[arrow] (if) -- (rules);

        % Entrée / Sortie
        \node[font=\footnotesize, align=center, left=0.4cm of io] (input) {\textbf{Données}\\[-7pt]\textbf{capteurs}};
        \node[font=\footnotesize, align=center, right=0.4cm of rules] (output) {\textbf{Épisodes}\\[-7pt]\textbf{\& notifs}};
        \draw[arrow] (input) -- (io);
        \draw[arrow] (rules) -- (output);
    \end{tikzpicture}
    \caption{Architecture du pipeline de détection de boiterie probable.}
    \label{fig:pipeline}
\end{figure}

Chaque étape transforme progressivement les données. Le chargement et filtrage normalise
les formats de fichiers et écarte les bins de couverture insuffisante. L'extraction de
\textit{features} convertit les signaux bruts en variables dérivées (z-scores robustes,
dérivées temporelles, moyennes mobiles) décrivant chaque vache relativement à son propre
comportement. La détection applique un Isolation Forest entraîné par vache et produit un
score d'anomalie par bin. Enfin, les règles métier (persistance temporelle, cohérence
multi-signaux, \textit{cooldown}) transforment ces anomalies en épisodes de boiterie probable et en
notifications. La suite de cette section détaille le prétraitement et l'extraction des
\textit{features}~; la détection et les règles métier font l'objet de la section suivante et
du Chapitre~\ref{chapitre:architecture}.

\FloatBarrier

Le prétraitement commence par une normalisation des noms de colonnes et des types temporels,
afin d'accepter les variantes de fichiers CSV réelles sans modifier le cœur du pipeline.
Les données sont ensuite ré-échantillonnées à l'intervalle de référence de 15~minutes,
qui correspond au pas temporel natif retenu pour la configuration finale. La construction
des variables dérivées vise à fournir à l'algorithme de détection une vue relative du
comportement de chaque vache, plutôt qu'une comparaison brute entre animaux.
Pour chaque variable brute $x$, le pipeline produit six familles de \textit{features}, détaillées
dans le Tableau~\ref{tab:features}.

\begin{table}[H]
    \centering
    \caption{Familles de \textit{features} dérivées pour chaque variable comportementale brute.}
    \label{tab:features}
    \begin{tabular}{lp{8cm}}
        \toprule
        \textbf{Suffixe}        & \textbf{Description} \\
        \midrule
        \texttt{\_rz}           & Z-score robuste global (MAD \citep{leys2013detecting}, replis IQR puis écart-type) \\
        \texttt{\_rrz}          & Z-score robuste glissant (fenêtre de 24 bins = 6\,h) \\
        \texttt{\_d1}           & Différence première (variation bin à bin) \\
        \texttt{\_d1\_per\_hour}& Vitesse de variation normalisée par heure \\
        \texttt{\_mm7}          & Moyenne mobile sur 7 bins ($\approx$1\,h\,45) \\
        \texttt{\_vs\_mm7}      & Écart à la moyenne mobile \\
        \bottomrule
    \end{tabular}
\end{table}

\noindent
S'y ajoutent des variables transversales:
\begin{itemize}
    \item \texttt{coverage\_pct}: taux de couverture capteur par bin, calculé à partir du
    nombre de lignes brutes observées rapporté au nombre attendu (médiane robuste);
    \item \texttt{sin\_hour}, \texttt{cos\_hour}, \texttt{sin\_dow}, \texttt{cos\_dow}:
    encodage cyclique de l'heure et du jour de semaine;
    \item \texttt{Motion Index\_sum\_log}, \texttt{Motion Index\_mean\_log}: transformations
    logarithmiques du \textit{Motion Index} pour limiter l'effet des valeurs extrêmes.
\end{itemize}

Au total, les six suffixes appliqués aux 7~variables brutes produisent $7 \times 6 = 42$
\textit{features} dérivées, auxquelles s'ajoutent les 7~variables transversales, soit environ
50~\textit{features} par bin. La Figure~\ref{fig:feature_engineering} illustre ce processus de
transformation pour une variable brute donnée. Chaque signal est ainsi converti en six
indicateurs complémentaires, la procédure étant répétée pour chacune des 7~variables
comportementales.

\begin{figure}[H]
    \centering
    \resizebox{0.88\textwidth}{!}{%
    \begin{tikzpicture}[
        box/.style={rectangle, draw=black!60, rounded corners=3pt, minimum height=0.85cm,
                    minimum width=2.8cm, align=center, font=\footnotesize, inner sep=12pt},
        arrow/.style={-{Stealth[length=2mm]}, thick, black!50},
        bus/.style={thick, black!50},
    ]
        % === Colonne 1 : Signal brut ===
        \node[box, fill=ifblue!12, minimum width=3cm] (brut)
            {Signal brut\\[-7pt]$x_t$ (ex: \textit{Steps})};
        \node[font=\scriptsize, black!45, below=4pt of brut] {$\times\,7$ variables};

        % === Colonne 2 : 6 transformations (pile verticale) ===
        \node[box, fill=loforange!12, right=3.2cm of brut, yshift=2.2cm] (d1h)
            {Vitesse/heure\\[-7pt]\texttt{\_d1\_per\_h}};
        \node[box, fill=loforange!12, below=0.35cm of d1h] (rrz)
            {Z-score glissant\\[-7pt]\texttt{\_rrz}};
        \node[box, fill=loforange!12, below=0.35cm of rrz] (rz)
            {Z-score global\\[-7pt]\texttt{\_rz}};
        \node[box, fill=loforange!12, below=0.35cm of rz] (d1)
            {Différence 1\textsuperscript{re}\\[-7pt]\texttt{\_d1}};
        \node[box, fill=loforange!12, below=0.35cm of d1] (mm7)
            {Moy. mobile 7j\\[-7pt]\texttt{\_mm7}};
        \node[box, fill=loforange!12, below=0.35cm of mm7] (vsmm)
            {Écart moy. mob.\\[-7pt]\texttt{\_vs\_mm7}};

        % Encadré arrière-plan (dessiné AVANT les noeuds via on background layer)
        \begin{scope}[on background layer]
            \node[draw=loforange!35, rounded corners=5pt, inner sep=10pt,
                  fit=(d1h)(vsmm), fill=loforange!4] (grp) {};
        \end{scope}
        \node[font=\scriptsize\bfseries, loforange!70!black, above=3pt of grp]
            {6 familles de \textit{features}};

        % === Colonne 3 : Vecteur final ===
        \node[box, fill=rulegreen!12, right=3.2cm of rz, minimum width=3cm,
              minimum height=2.4cm] (vect)
            {Vecteur\\[-5pt]de \textit{features}\\[-5pt]($\approx$\,50 dim.)\\[-5pt]par bin};

        % === Bus gauche (distribution) : ligne verticale unique ===
        \coordinate (busL) at ([xshift=1.0cm]brut.east);
        \draw[bus] (brut.east) -- (busL);
        \draw[bus] (busL |- d1h.west) -- (busL |- vsmm.west);
        \foreach \nd in {d1h,rrz,rz,d1,mm7,vsmm}
            \draw[arrow] (busL |- \nd.west) -- (\nd.west);

        % === Bus droit (collecte) : ligne verticale unique ===
        \coordinate (busR) at ([xshift=-1.0cm]vect.west);
        \draw[bus] (busR |- d1h.east) -- (busR |- vsmm.east);
        \foreach \nd in {d1h,rrz,rz,d1,mm7,vsmm}
            \draw[arrow] (\nd.east) -- (busR |- \nd.east);
        \draw[arrow, thick] (busR) -- (vect.west);

    \end{tikzpicture}%
    }
    \caption{Processus d'extraction de \textit{features} pour une variable brute.}
    \label{fig:feature_engineering}
\end{figure}

L'ingénierie de variables suit une logique pragmatique qui consiste à transformer des signaux
capteurs imparfaits en indicateurs comportementaux plus stables. Des travaux méthodologiques sur le
couchage montrent par exemple que l'intervalle d'échantillonnage et les règles d'édition
modifient déjà la qualité des mesures dérivées \citep{ledgerwood2010lying}. À l'inverse, des
travaux plus récents sur accélérométrie tri-axiale montrent qu'une modélisation plus riche
devient surtout utile lorsqu'on vise une classification fine de comportements plutôt qu'une
simple alerte de déviation \citep{balasso2023behaviour}. Dans ce projet, la priorité est la
robustesse opérationnelle et l'interprétabilité d'une alerte probable, pas une taxonomie
complète des comportements.

L'objectif n'est pas de multiplier artificiellement les variables, mais de capturer trois
dimensions complémentaires: le niveau d'activité, la rupture par rapport au comportement
habituel, et la structure temporelle locale. Cette représentation est adaptée à une
logique de détection d'anomalies intra-animal, plus défendable qu'une comparaison
inter-animal brute dans un contexte sans étiquettes cliniques exhaustives.

\FloatBarrier
% =====================================================================
\section{Détection et règles métier}
% =====================================================================

Le cœur algorithmique repose sur Isolation Forest (IF), choisi comme détecteur
d'anomalies non supervisé adapté à un contexte de labels partiels ou absents
\citep{liu2008isolation}. Le modèle est entraîné individuellement par vache sur une
baseline initiale (\texttt{baseline\_ratio} = 0.6). Le score IF est ensuite filtré par
des règles métier (persistance temporelle, cohérence multi-signaux, espacement des
notifications) dont le détail algorithmique est présenté au
Chapitre~\ref{chapitre:architecture}. La configuration finale figée dans le projet
est détaillée dans le Tableau~\ref{tab:config_finale}.

\begin{table}[H]
    \centering
    \caption{Configuration finale du pipeline, gelée après validation.}
    \label{tab:config_finale}
    \begin{tabular}{llp{6.5cm}}
        \toprule
        \textbf{Paramètre}           & \textbf{Valeur} & \textbf{Rôle} \\
        \midrule
        \texttt{interval}            & 15T             & Pas temporel de ré-échantillonnage \\
        \texttt{contamination}       & 0.06            & Fraction d'anomalies attendue par IF \\
        \texttt{baseline\_ratio}     & 0.6             & Part de la série utilisée comme baseline \\
        \texttt{persist\_hours}      & 7               & Fenêtre de persistance (heures) \\
        \texttt{alert\_min}          & 2               & Bins anomaux minimaux par fenêtre \\
        \texttt{mix\_mode}           & MIX             & Mode de cohérence multi-signaux \\
        \texttt{mix\_rate\_thr}      & 0.24            & Seuil de taux d'anomalies \\
        \texttt{mi\_z\_high\_thr}    & 2.2             & Seuil z-score \textit{Motion Index} \\
        \texttt{cooldown\_hours}     & 12              & Espacement minimal entre notifications \\
        \texttt{coverage\_min\_pct}  & 25              & Couverture minimale pour bin fiable (\%) \\
        \texttt{random\_state}       & 42              & Seed pour reproductibilité IF \\
        \bottomrule
    \end{tabular}
\end{table}

Les seuils retenus n'ont pas été fixés arbitrairement; ils résultent d'une exploration progressive
des configurations via les validations v4, v5.1 et v5.2 (Chapitre~\ref{chapitre:validation}),
puis d'un gel explicite garantissant la reproductibilité.

La Figure~\ref{fig:persistance_cooldown} illustre le fonctionnement conjoint de la
persistance temporelle et du cooldown sur un exemple temporel concret. Les anomalies IF
(en rouge) ne déclenchent une notification que si elles sont assez nombreuses dans la
fenêtre de persistance (7~h)~; le cooldown de 12~h empêche ensuite les alertes
répétitives, si bien que l'anomalie isolée (h15) et le groupe h18--19 (cooldown actif)
ne déclenchent rien. Cette mécanique est essentielle pour comprendre pourquoi le
pipeline ne réagit pas aux anomalies isolées et pourquoi les notifications sont espacées
dans le temps.

\begin{figure}[H]
    \centering
    \begin{tikzpicture}[x=0.47cm, y=0.6cm]
        % Axe temps
        \draw[-{Stealth[length=2.5mm]}, black!70] (0,0) -- (28,0);
        \node[font=\scriptsize, black] at (28,-0.7) {Temps (heures)};
        \foreach \i in {0,4,8,12,16,20,24} {
            \draw[black!40] (\i,0) -- (\i,-0.15);
            \node[font=\scriptsize, black!50, below] at (\i,-0.15) {\i};
        }

        % Anomalies IF (bins individuels)
        \node[font=\scriptsize\bfseries, black, anchor=east] at (-0.3,1.2) {Anomalies IF};
        \foreach \i in {2,3,5,6,7,8,9,10,11,18,19} {
            \fill[rawred!50] (\i,0.7) rectangle (\i+0.8,1.7);
        }
        % Anomalie isolée
        \fill[profgray!40] (15,0.7) rectangle (15.8,1.7);
        \node[font=\tiny, black!50, above=3pt] at (15.4,1.7) {isolée};

        % Fenêtre de persistance
        \node[font=\scriptsize\bfseries, black, anchor=east] at (-0.3,3.2) {Persistance};
        \draw[ifblue, thick] (2,2.7) rectangle (9,3.7);
        \node[font=\scriptsize, ifblue!80!black] at (5.5,3.2) {7h};
        \draw[decorate, decoration={brace, amplitude=4pt, raise=2pt}, ifblue!70]
            (2,3.7) -- (9,3.7)
            node[midway, above=7pt, font=\tiny, ifblue!70!black] {seuil atteint};

        % Notifications
        \node[font=\scriptsize\bfseries, black, anchor=east] at (-0.3,5.4) {Notifications};
        \fill[rulegreen!30] (9,4.9) rectangle (9.8,5.9);
        \draw[rulegreen, thick] (9,4.9) rectangle (9.8,5.9);
        \node[font=\tiny\bfseries, rulegreen!80!black, above=3pt] at (9.4,5.9) {N1};

        % Cooldown
        \draw[loforange, thick, densely dashed] (9.8,5.4) -- (21.8,5.4);
        \draw[decorate, decoration={brace, amplitude=4pt, mirror, raise=8pt}, loforange!70]
            (9.8,4.9) -- (21.8,4.9)
            node[midway, below=12pt, font=\scriptsize, loforange!80!black] {Cooldown (12h)};

        % Deuxième notification bloquée
        \node[font=\tiny, rawred!70, above=3pt] at (19,5.9) {bloquée};
        \draw[rawred!50, thick] (18.6,5.0) -- (19.4,5.8);
        \draw[rawred!50, thick] (18.6,5.8) -- (19.4,5.0);

    \end{tikzpicture}
    \caption{Mécanisme de persistance et de cooldown.}
    \label{fig:persistance_cooldown}
\end{figure}

\subsection{Normalisation robuste}
La normalisation des \textit{features} avant détection utilise le \texttt{RobustScaler} de
scikit-learn \citep{sklearn_api}, qui centre les données sur la médiane et les normalise
par l'écart interquartile (IQR). Ce choix est motivé par la présence de valeurs extrêmes
dans les signaux capteurs (cf. Figure~\ref{fig:distribution_steps}): contrairement au
\texttt{StandardScaler} (moyenne et écart-type), le \texttt{RobustScaler} n'est pas
influencé par les outliers \citep{leys2013detecting}, ce qui évite de comprimer
artificiellement la distribution des \textit{features} normales. Les quantiles 10--90 sont utilisés
plutôt que les quartiles classiques 25--75, afin de capturer une plus grande partie de
la distribution tout en restant robuste aux extrêmes.

\subsection{Stratégie d'entraînement par vache}
Le modèle IF est entraîné individuellement pour chaque vache, en utilisant uniquement les
60\% premiers bins fiables comme baseline (\texttt{baseline\_ratio} = 0.6). Cette stratégie
présente deux avantages. D'une part, elle capture le comportement normal
spécifique de chaque animal, sans dilution par les profils des autres vaches. D'autre part,
elle isole la phase d'entraînement de la phase de prédiction: le modèle est calibré sur
une période supposée normale, puis appliqué sur le reste de la série. Ce schéma est cohérent
avec l'hypothèse que la boiterie est un événement rare et tardif, plutôt qu'un état présent
dès le début de l'observation.

Enfin, une variante LOF (Local Outlier Factor) a été évaluée dans l'étude
d'ablation comme point de comparaison méthodologique \citep{breunig2000lof}. Elle n'est
pas la baseline officielle du projet, mais sert à tester la sensibilité des résultats au
choix du détecteur d'anomalies.

\FloatBarrier
% =====================================================================
\section{Métriques d'évaluation}
% =====================================================================

L'évaluation a été construite pour rester honnête vis-à-vis des données disponibles. Cette
prudence est conforme aux recommandations de la littérature en détection automatisée de
boiterie, qui souligne les limites des métriques classiques (sensibilité, spécificité)
lorsque le gold standard clinique n'est pas disponible \citep{alsaaod2019automatic,
rutten2013invited}. En l'absence de ce gold standard, le projet emploie des métriques
événementielles et opérationnelles plutôt que des métriques cliniques au sens fort,
résumées dans le Tableau~\ref{tab:metriques}.

\begin{table}[H]
    \centering
    \caption{Métriques d'évaluation utilisées dans les protocoles de validation.}
    \label{tab:metriques}
    \small
    \begin{tabular}{>{\ttfamily\raggedright\arraybackslash}p{3.8cm}
                    >{\rmfamily\raggedright\arraybackslash}p{8cm}}
        \toprule
        \textrm{\textbf{Métrique}}             & \textbf{Définition} \\
        \midrule
        detected\_any\_overlap      & L'épisode injecté chevauche au moins un
                                               épisode prédit (binaire). \\
        detected\_iou20             & Le chevauchement atteint un seuil strict
                                               de 20\,\% (IoU $\geq$ 0.20). \\
        best\_iou                   & Meilleur recouvrement continu (Intersection
                                               over Union) obtenu sur l'événement. \\
        false\_notif\_cow\_day      & Nombre de notifications hors fenêtre injectée,
                                               normalisé par vache-jour. \\
        hidden\_leak                & Fuite sur scénario volontairement non
                                               détectable (sensibilité au bruit). \\
        \bottomrule
    \end{tabular}
\end{table}

Un tel choix met l'accent sur un compromis réaliste: détecter les cas forts, conserver une
sensibilité raisonnable sur les cas limites, et rester prudent sur les fausses alertes.

\subsection{Intersection over Union (IoU)}
L'Intersection over Union (IoU) est la métrique centrale de recouvrement temporel, empruntée
à la détection d'objets en vision par ordinateur \citep{everingham2010pascal}. Elle mesure le
degré de chevauchement entre un épisode injecté et un épisode prédit, en rapportant
l'intersection des deux fenêtres à leur union:

\begin{equation}
    \text{IoU} = \frac{|A \cap B|}{|A \cup B|}
    = \frac{\text{bins communs}}{\text{bins injectés} + \text{bins prédits} - \text{bins communs}}
\end{equation}

\noindent
où $A$ est la fenêtre de l'épisode injecté et $B$ la fenêtre de l'épisode prédit (en
nombre de bins de 15~minutes). Un IoU de 1 indique un recouvrement parfait; un IoU de 0
indique une absence totale de chevauchement. Le seuil \texttt{IoU20} retenu dans ce projet
($\text{IoU} \geq 0.20$) signifie qu'au moins 20\% de la fenêtre combinée est couverte
par la prédiction. Ce seuil a été choisi pour capturer les détections partielles mais
significatives, tout en rejetant les chevauchements marginaux.

La Figure~\ref{fig:iou_concept} illustre ce concept sur un exemple temporel~: avec
3~bins d'intersection sur 10~bins d'union, l'IoU vaut 0.30, ce qui dépasse le seuil de
20\% et valide la détection au sens strict.

\begin{figure}[H]
    \centering
    \begin{tikzpicture}[x=0.8cm, y=0.8cm]

        % --- Grille de fond (bins) ---
        \foreach \i in {0,1,...,15} {
            \draw[black!8] (\i,0) -- (\i,4.5);
        }

        % --- Graduation axe temps ---
        \draw[-{Stealth[length=2.5mm]}, thick, black!70] (-0.3,0) -- (15.5,0);
        \foreach \i in {0,3,6,9,13} {
            \draw[black!50] (\i,0) -- (\i,-0.15);
            \node[font=\scriptsize, black!60, below] at (\i,-0.15) {\i};
        }
        \node[font=\scriptsize, black] at (15.5,-0.5) {Temps (bins)};

        % --- Episode injecte (A) ---
        \fill[ifblue!18] (3,2.8) rectangle (9,3.8);
        \draw[ifblue, thick] (3,2.8) rectangle (9,3.8);
        \node[font=\small, black] at (6,3.3) {Épisode injecté ($A$)};
        % Accolade
        \draw[decorate, decoration={brace, amplitude=5pt, raise=2pt}, black!70]
            (3,3.8) -- (9,3.8) node[midway, above=7pt, font=\scriptsize, black] {6 bins};

        % --- Episode predit (B) ---
        \fill[rulegreen!18] (6,0.8) rectangle (13,1.8);
        \draw[rulegreen, thick] (6,0.8) rectangle (13,1.8);
        \node[font=\small, black] at (9.5,1.3) {Épisode prédit ($B$)};
        % Accolade
        \draw[decorate, decoration={brace, amplitude=5pt, mirror, raise=2pt}, black!70]
            (6,0.8) -- (13,0.8) node[midway, below=7pt, font=\scriptsize, black] {7 bins};

        % --- Zone d'intersection (A inter B) ---
        \fill[loforange!25] (6,1.8) rectangle (9,2.8);
        \draw[loforange, thick] (6,1.8) rectangle (9,2.8);
        \node[font=\small\bfseries, black] at (7.5,2.3) {$A \cap B$};
        % Accolade laterale
        \draw[decorate, decoration={brace, amplitude=4pt, raise=2pt}, black!70]
            (9,2.8) -- (9,1.8) node[midway, right=6pt, font=\scriptsize, black] {3 bins};

        % --- Lignes de liaison verticales ---
        \draw[black!25, dashed] (6,0.8) -- (6,3.8);
        \draw[black!25, dashed] (9,0.8) -- (9,3.8);

        % --- Formule de calcul ---
        \node[anchor=west, font=\small, black] at (0,-1.2)
            {$\displaystyle \text{IoU} = \frac{|A \cap B|}{|A \cup B|}
            = \frac{3}{6 + 7 - 3} = \frac{3}{10} = 0.30$};

        % --- Verdict ---
        \node[anchor=west, font=\small\bfseries, black] at (0,-2.3)
            {$0.30 \geq 0.20 \quad \Rightarrow \quad$ IoU20 = 1 \enspace (détection validée)};

    \end{tikzpicture}
    \caption{Illustration du calcul d'IoU entre un épisode injecté et un épisode prédit.}
    \label{fig:iou_concept}
\end{figure}

\FloatBarrier
% =====================================================================
\section{Protocole d'injection}
% =====================================================================

La validation repose sur des injections synthétiques dans les données capteurs réelles. Cette
approche, analogue aux protocoles de contamination artificielle utilisés en évaluation de
détecteurs d'anomalies \citep{emmott2015metaanalysis,chandola2009anomaly}, permet de
contrôler la nature, l'amplitude et la durée des perturbations, et donc de mesurer les
performances du pipeline dans des conditions connues. Trois profils d'injection ont été
conçus pour couvrir un spectre de difficulté (Tableau~\ref{tab:injection_profiles}):

\begin{table}[H]
    \centering
    \caption{Profils d'injection utilisés dans les validations manuelles et robustes.}
    \label{tab:injection_profiles}
    \small
    \begin{tabular}{>{\ttfamily\raggedright\arraybackslash}p{4.1cm}
                    >{\raggedright\arraybackslash}p{3.7cm}cc}
        \toprule
        \textbf{\textrm{Profil}}     & \textbf{Description}
            & \textbf{Durée} & \textbf{Dét. att.} \\
        \midrule
        detectable\_strong  & Épisode fort, multi-signaux,
                                       amplitude élevée (2--3$\times$ la normale)
            & $2.0 \times K$ & Oui \\
        detectable\_borderline & Épisode faible, amplitude modérée
                                          (1.3--2$\times$), plus court
            & $1.45 \times K$ & Partielle \\
        non\_detectable\_short & Perturbation brève, mono-signal,
                                          proche du bruit
            & $0.3 \times K$ & Non \\
        \bottomrule
    \end{tabular}
\end{table}

\noindent
où $K = \lceil \texttt{persist\_hours} \times 60 / \texttt{interval\_min} \rceil$ est le
nombre de bins dans la fenêtre de persistance. Le profil \texttt{strong} modifie cinq
familles de signaux en alternant phases normales et phases élevées, simulant un épisode
multi-signaux cohérent. Le profil \texttt{borderline} utilise des amplitudes plus faibles,
proches du seuil de détection. Le profil \texttt{non\_detectable\_short} ne modifie qu'un
seul signal sur une durée trop courte pour déclencher la logique de persistance.

Une telle gradation est essentielle: elle mesure non seulement la sensibilité du
système (cas forts), mais aussi sa spécificité (cas non détectables) et sa zone grise
(cas borderline). Un système qui détecte tout, y compris les cas non détectables, serait
trop sensible et inutilisable en pratique.

\vspace{0.5em}
\noindent
Les données, les \textit{features} et le protocole expérimental étant définis, le chapitre suivant
décrit l'architecture logicielle du pipeline, les choix d'implémentation et les détails
algorithmiques des deux détecteurs retenus (Isolation Forest et Local Outlier Factor).
L'objectif est de montrer comment les décisions méthodologiques présentées ici se traduisent
en composants modulaires, testables et reproductibles.
